% Generated draft: requires mathematical and source review.
\documentclass[11pt]{article}
\usepackage[margin=25mm]{geometry}
\usepackage{fontspec}
\setmainfont{Latin Modern Roman}
\usepackage{amsmath,amssymb,mathtools,unicode-math}
\setmathfont{Latin Modern Math}
\usepackage{enumitem,needspace,xurl,hyperref}
\setlength{\parindent}{0pt}
\setlength{\parskip}{5pt}
\setlength{\emergencystretch}{3em}
\newcommand{\N}{\mathbb{N}}
\newcommand{\Z}{\mathbb{Z}}
\newcommand{\Q}{\mathbb{Q}}
\newcommand{\R}{\mathbb{R}}
\newcommand{\C}{\mathbb{C}}
\newcommand{\F}{\mathbb{F}}
\newcommand{\A}{\mathbb{A}}
\newcommand{\PP}{\mathbb{P}}
\newcommand{\E}{\mathbb{E}}
\newcommand{\eps}{\varepsilon}
\newcommand{\dd}{\,\mathrm{d}}
\begin{document}
\section*{UnsolvedMath ID: 20003286}
\textbf{Source problem code:} catalogue entry
\subsection{Definitions and mathematical statement}
This is a status paragraph following Question 17, not a second independent question. It reports a result of N. Shah for $L=SO(m,1)$ and $G=SO(n,1)$, $m>n$, under the preceding homogeneous-dynamics setup. Neighboring local record $20003285$ identifies a lattice quotient of $L$, a semisimple $a\in G$, its expanding horospherical subgroup $U$, and expanding curve measures $a^j\nu$. The present fragment says these converge to Haar measure $\lambda$ if the curve image lies in neither a proper affine subspace nor an $(n-2)$-sphere in $U\cong\R^{n-1}$. The basepoint, orbit-density, component and measure hypotheses of the full Shah result are not supplied here, so it is not asserted for arbitrary quotients and curves without those inputs. The iteration index is written $j$ to distinguish it from dimension $n$.

The aid's exact certificate concerns the stated geometric condition alone. Put $d=n-1\geq1$ and let $\gamma:I\to\R^d$, with $I$ an interval. Define its quadratic lift
\[
 v(t)=(1,\gamma_1(t),\ldots,\gamma_d(t),\|\gamma(t)\|^2)
           \in\R^{d+2}.
\]
The image is contained in neither any proper affine subspace nor any round $(d-1)$-sphere precisely when the functions
\[
 1,\gamma_1,\ldots,\gamma_d,\|\gamma\|^2
\]
are linearly independent over $\R$. A relation $c+b\cdot\gamma+a\|\gamma\|^2=0$ with $a=0$ gives a proper affine hyperplane. With $a\ne0$, completing the square gives a sphere; a zero radius is a point, already affine-degenerate, and a negative radius-square cannot contain a nonempty curve image.

Equivalently, there exist $d+2=n+1$ parameter values $t_0,\ldots,t_{d+1}$ for which the matrix with rows $v(t_i)$ has nonzero determinant. This is an existence certificate, not a guarantee for every preassigned sample set. The number is worst-case minimal for a value-only test: any at most $d+1$ sample points lie in a proper affine subspace or, if affinely spanning, on a sphere, and cannot certify full rank of the lift.

If $\gamma$ is real analytic on a connected open interval, define $f=(1,\gamma_1,\ldots,\gamma_d,\|\gamma\|^2)$. The same condition is equivalent to the standard Wronskian
\[
 W(t)=\det\bigl(f_j^{(r)}(t)\bigr)_{0\leq r,j\leq d+1}
\]
being not identically zero. It may still vanish at particular interior parameters. At every prescribed interior $t_*$, the exact replacement is existence of derivative orders
\[
 0=r_0<r_1<\cdots<r_{d+1}
\]
with $\det(f_j^{(r_i)}(t_*))\ne0$. Analyticity ensures that a linear combination with all derivatives zero at $t_*$ vanishes identically, which is the needed full-jet principle. No uniform bound on these orders is supplied. These are certificates for the curve hypothesis, not an independent proof of measure equidistribution or an extension to merely smooth curves.
\subsection{Short English statement}
The entry reports Shah's rank-one answer to a preceding equidistribution question. The attached aid gives exact tests for its geometric curve condition: independence of the quadratic-lift functions, a nonzero determinant at some $n+1$ samples, or analytic derivative certificates. A standard Wronskian can vanish at one point even when the curve is nondegenerate.
\subsection{Sources}
\begin{itemize}
\item Local UnsolvedMath registry, permanent ID $20003286$, source code catalogue entry; canonical location \texttt{aim-other-notes.json}, \texttt{notes[92]}. Neighboring record $20003285$ supplies part of Question 17's homogeneous setup.
\item AIM \emph{Emerging applications of measure rigidity}, section 4.3, explanatory paragraph immediately following Question 17, \url{https://aimath.org/WWN/measrigid/measrigid.pdf}. It attributes the stated rank-one result to N. Shah. No full external citation is supplied.
\item Research attempt 1 is marked context-only/not-a-problem, proved, and low novelty confidence, despite the registry question/open labels. Its explicit caveat says the aid is machine generated, not peer reviewed, and requires independent verification before citation.
\end{itemize}
\end{document}
